% !TeX program = xelatex
\documentclass[UTF8,12pt,a4paper]{ctexart}
\usepackage{amsmath,amssymb}
\usepackage{geometry}
\geometry{margin=2cm}

% ===== 水印：每页3个位置（右上/中间/右下）=====
\usepackage{xcolor}
\usepackage{eso-pic}
\usepackage{graphicx}

\newcommand{\WMText}{贺昱琪学习资料}
\newcommand{\WMScale}{2.28}
\newcommand{\WMAngle}{45}
\colorlet{WMColor}{black!8}

\AddToShipoutPictureBG{%
  \AtPageLowerLeft{%
    \put(\LenToUnit{0.66\paperwidth},\LenToUnit{0.70\paperheight}){%
      \rotatebox{\WMAngle}{\scalebox{\WMScale}{\textcolor{WMColor}{\WMText}}}%
    }%
    \put(\LenToUnit{0.33\paperwidth},\LenToUnit{0.50\paperheight}){%
      \rotatebox{\WMAngle}{\scalebox{\WMScale}{\textcolor{WMColor}{\WMText}}}%
    }%
    \put(\LenToUnit{0.66\paperwidth},\LenToUnit{0.30\paperheight}){%
      \rotatebox{\WMAngle}{\scalebox{\WMScale}{\textcolor{WMColor}{\WMText}}}%
    }%
  }%
}

% ===== 插图兜底：图片不存在也不报错 =====
\newcommand{\SafeImg}[2]{%
  \IfFileExists{#1}{%
    \begin{center}
      \vspace{-0.4em}
      \includegraphics[width=#2\linewidth]{#1}
      \vspace{-0.6em}
    \end{center}
  }{%
    \begin{center}
      \vspace{-0.2em}
      \fbox{\textit{缺少图片：#1（请将图片与.tex放同一文件夹，并命名一致）}}
      \vspace{-0.2em}
    \end{center}
  }%
}

% ===== 5题铺满一页 + 底部留白 =====
\usepackage{enumitem}
\newcommand{\BottomWriteSpace}{1.8cm}

\newenvironment{OnePageFiveFill}{
  \small
  \vspace*{0pt}
  \begin{enumerate}[leftmargin=*, itemsep=0pt, topsep=0pt, parsep=0pt, partopsep=0pt]
  \setlength{\parskip}{0pt}
}{
  \end{enumerate}
  \vspace*{\BottomWriteSpace}
  \normalsize
  \newpage
}

\newcommand{\Q}[1]{\item #1\par\vfill}
\newcommand{\Qlast}[1]{\item #1\par}

% ===== 特殊留白样式 =====
\newcommand{\BlankSpace}{\vspace{2.6cm}}

\begin{document}

% =========================================================
% 4/7
% =========================================================
\section*{4月7日作业（5道题当晚22:00前打卡）}
\begin{OnePageFiveFill}

% 2道原难度题目
\Q{$\displaystyle \frac{1}{1\cdot3} + \frac{1}{3\cdot5} + \frac{1}{5\cdot7} + \cdots + \frac{1}{97\cdot99}$}

\Q{$\dfrac{2x-1}{3} - \dfrac{5x+1}{6} = 1$}

% 3道小升初基础计算
\Q{$4.5 \times 102$}

\Q{$\displaystyle \frac{5}{8} \times \frac{4}{15} \div \frac{1}{6}$}

\Qlast{$0.125 \times 0.25 \times 320$}

\end{OnePageFiveFill}

% =========================================================
% 4/8
% =========================================================
\section*{4月8日作业（5道题当晚22:00前打卡）}
\begin{OnePageFiveFill}

% 2道原难度题目
\Q{$\displaystyle \left(1+\frac{1}{2}\right) \times \left(1+\frac{1}{3}\right) \times \left(1+\frac{1}{4}\right) \times \cdots \times \left(1+\frac{1}{99}\right)$}

\Q{$0.3(4x-1) = 0.5(x+5) - 0.8$}

% 3道小升初基础计算
\Q{$9.8 \times 25$}

\Q{$\displaystyle \frac{4}{9} + \frac{1}{6} + \frac{5}{9} + \frac{5}{6}$}

\Qlast{$3.6 \div 0.12$}

\end{OnePageFiveFill}

% =========================================================
% 4/9
% =========================================================
\section*{4月9日作业（5道题当晚22:00前打卡）}
\begin{OnePageFiveFill}

% 2道原难度题目
\Q{$\displaystyle \frac{2}{2\cdot4} + \frac{2}{4\cdot6} + \frac{2}{6\cdot8} + \cdots + \frac{2}{48\cdot50}$}

\Q{$(1.5 \times 0.4 + 0.6 \times 3) : x = 4 : 5$}

% 3道小升初基础计算
\Q{$\displaystyle 24 \times \left(\frac{1}{4} - \frac{1}{6}\right)$}

\Q{$7.5 \times 4.8 + 5.2 \times 7.5$}

\Qlast{$\displaystyle 1 \div \left( \frac{1}{2} \div \frac{1}{3} \right)$}

\end{OnePageFiveFill}

% =========================================================
% 4/10
% =========================================================
\section*{4月10日作业（5道题当晚22:00前打卡）}
\begin{OnePageFiveFill}

% 2道原难度题目
\Q{$\displaystyle \frac{1}{2} + \frac{1}{6} + \frac{1}{12} + \frac{1}{20} + \frac{1}{30} + \frac{1}{42} + \frac{1}{56}$}

\Q{$\dfrac{x}{2} + \dfrac{x}{3} + \dfrac{x}{4} = 13$}

% 3道小升初基础计算
\Q{$101 \times 0.99$}

\Q{$\displaystyle \frac{7}{15} \times \frac{5}{14} + \frac{1}{6}$}

\Qlast{$8.4 \div (0.2 \times 0.5)$}

\end{OnePageFiveFill}

% =========================================================
% 4/11
% =========================================================
\section*{4月11日作业（5道题当晚22:00前打卡）}
\begin{OnePageFiveFill}

% 2道原难度题目
\Q{$\displaystyle \frac{2022 \times 2023 - 1}{2022 \times 2022 + 2021}$}

\Q{$\dfrac{3x-2}{4} - 1 = \dfrac{5x+3}{8}$}

% 3道小升初基础计算
\Q{$5.6 \times 1.25$}

\Q{$\displaystyle \left( \frac{3}{8} + \frac{1}{4} \right) \div \frac{5}{16}$}

\Qlast{$99.9 \times 9 + 9.99 \times 10$}

\end{OnePageFiveFill}

\end{document}